\documentclass[10pt]{article}
\usepackage[margin=0.8in]{geometry}
\usepackage{amsmath,booktabs,xcolor,enumitem,graphicx}
\usepackage[hidelinks]{hyperref}
\setlist{nosep,leftmargin=1.15em}
\pagestyle{empty}

\begin{document}
\begin{center}
{\LARGE\textbf{RMFS at a glance}}\\[2pt]
{\small Robust Multilingual Factor Score --- Waiz Khan, JHU,
August 2026}
\end{center}
\vspace{-4pt}\hrule\vspace{6pt}

\textbf{What it is.} A report card for a model's multilingual
representations: four readings per language, each with its own
validated number, scored 0 to 100 against a reference grid of 14
models. A score of 74 means better than 74 percent of the reference
models.

\textbf{How it is computed} (about four GPU minutes per model): embed
300 parallel sentences in 41 languages at one layer, then take:
\begin{itemize}
\item \textbf{Meaning against language} --- the share of representation
structure that encodes concepts rather than language identity.
\item \textbf{Sentence matching} --- how reliably the model matches
sentences to their English parallels.
\item \textbf{Weakest aligned languages} --- behaviour in the poorest
aligned tail, where training damage appears first.
\item \textbf{Collapse flag} --- meaning structure against a reference
checkpoint, for training settings.
\end{itemize}

\textbf{What each reading is worth.} Correlations with outside
evidence, after removing training data volume, language family, script
and tokenizer effects. No reading wins twice, which is why they are
reported side by side rather than averaged into one score:
\begin{center}\small
\begin{tabular}{lccc}
\toprule
 & carrying content & downstream & trained variants\\
 & across languages & exams & (75 checkpoints)\\
\midrule
Meaning against language & $\mathbf{+0.76}$ & $+0.07$ & $-0.10$\\
Sentence matching & $+0.76$ & $\mathbf{+0.24}$ & $+0.36$\\
Weakest languages & --- & $+0.09$ & $\mathbf{+0.51}$\\
Collapse flag & --- & --- & detects exactly\\
\bottomrule
\end{tabular}
\end{center}

\textbf{Scoreboard.} Qwen3-8B \textbf{90.5}, 4B \textbf{85.7}, 1.7B
\textbf{73.8}; Salamandra-7B \textbf{69.0}; bloom-7b1 and Qwen3-0.6B
\textbf{54.8}; OLMo-2-1B \textbf{28.6}. Scale and multilingual design
orderings come out with no tuning.

\textbf{Why four readings and not one number.} In a study of 25
training variants (75 checkpoints), the variant trained on the word
level alignment objective lost about 95 percent of its meaning
structure and posted the best single score in the study; MEXA ranked it
third of 25. A share cannot see a loss that shrinks both of its parts,
and the same holds for any scale invariant ratio. That variant also
kept its downstream ability, 0.395 against 0.390 for the untouched
control, while the variants that did lose ability looked clean on every
geometric measure.

\textbf{What the numbers rest on.} Downstream prediction sits near 0.3
after controls for every measure in this area. The 0.76 holds when
whole language families are resampled (0.60 to 0.89) and rests on four
models in that comparison, giving 0.12 to 0.84 when models are the
unit. A panel of 69 languages carries the independent evidence of
roughly two.

\end{document}
